\begin{definition}
Let $n\ge2$, $K\ge1$, let $\Omega\subseteq\R^n$ be a domain, let
$F\colon\R^n\to\R^n$, let $G^1,\dots,G^n$ be the component weak gradients of $F$
and let $\tau\colon\R^n\to\R^{n\times n}$.  The predicate
$\mathrm{CitedQuasiregularBasics}(n,K,\Omega,F,G,\tau)$ is the cited result (C7)
of the Note for the formalization ([BI], [R], [GLM]):

\emph{if $F$ is $K$-quasiregular on $\Omega$ (\noderef{IsQuasiregularOn}) with
distortion matrix $\tau$ (\noderef{IsDistortionMatrix}), then}
\begin{enumerate}
\item \emph{$F$ agrees almost everywhere on $\Omega$ with a function that is
  continuous on $\Omega$: there is $\widetilde F$ with $\widetilde F$ continuous
  on $\Omega$ and $F=\widetilde F$ almost everywhere on $\Omega$;}
\item \emph{if $F$ is not almost everywhere constant on $\Omega$ --- that is, if
  there is no $c$ with $F=c$ almost everywhere on $\Omega$ --- then $J_F>0$
  almost everywhere on $\Omega$;}
\item \emph{each component $F^i$ is a weak solution in $\Omega$ of
\[
\operatorname{div}\Bigl((\tau(x)\nabla u,\nabla u)^{\frac n2-1}\tau(x)\nabla u\Bigr)=0
\tag{5.11}
\]
(\noderef{IsWeakSolutionOn}); and}
\item \emph{each component $F^i$ is a local minimum in $\Omega$ of
\[
v\longmapsto\int_E(\tau(x)\nabla v,\nabla v)^{\frac n2}\,dx
\]
(\noderef{IsLocalMinimumOn}).}
\end{enumerate}
Here $J_F(x)=\det DF(x)$ with $DF(x)_{ij}=G^i(x)^j$, and
$(\tau(x)h,h)=\sum_{a,b}\tau(x)_{ab}h^ah^b$.  The exponent $n/2$ in the
functional of (4) is the corrected one recorded in the source (the print has
$n$; the Euler equation of the functional is (5.11) of the thesis, and the functional
must have the growth (F1) of the thesis of \S4 with $p=n$).

\medskip
\noindent\textbf{Why clauses (1) and (2) are stated almost everywhere
invariantly.}  (C7), at paper l.~268, reads ``A $K$-quasiregular mapping $F$ is
continuous; if $F$ is not constant then $J_F>0$ a.e.''  Here $F$ is the
$W^{1,n}_{\mathrm{loc}}$ map of \noderef{IsQuasiregularOn}, and a
$W^{1,n}_{\mathrm{loc}}$ map is determined only up to a null set: every other
clause of this predicate --- the weak-gradient identities of clause (3), the
integral inequality of clause (4), the membership clauses inside
\noderef{IsQuasiregularOn} --- is unchanged if $F$ is redefined on a set of
measure zero.  Asserting \emph{of the given representative} that it is
continuous, or that it is not constant at every point, is therefore not an
assertion about the quasiregular map at all; and in conclusion position it makes
the predicate unsatisfiable, since from any $F$ for which the other clauses hold
one may change the value at a single point and destroy continuity while keeping
everything else.  The faithful reading of ``$F$ is continuous'' for an
almost-everywhere class is ``$F$ has a continuous representative'', which is
clause (1) above; and the matching reading of ``$F$ is not constant'' is that no
constant agrees with $F$ almost everywhere, which is clause (2)'s antecedent.
The Note at l.~258 requires the cited results to be stated as in the literature
and not in a stronger form, and the pointwise readings are strictly stronger:
they silently restrict Theorem~Theorem~7 of the thesis to those representatives that happen
to be continuous.  The same interface decision was taken for boundary values in
\noderef{MemW1pZeroOn}'s setting and is the shape
\noderef{CitedCampanatoEmbedding} already uses, where (C5)'s own wording at
l.~266 is ``agrees a.e.\ with a function in $C^{0,\alpha}$''.

\noindent\textbf{On the junk-value guard for clauses (3) and (4).}  Those two
clauses are the only place in the development where
\noderef{IsWeakSolutionOn} and \noderef{IsLocalMinimumOn} are instantiated
without an (F1) of the thesis or (Q1) of the thesis--(Q5) of the thesis bundle beside them: the
density here is $(\tau(x)h,h)^{n/2}$ and the field of (5.11) of the thesis is built
from $\tau$, which \noderef{IsDistortionMatrix} constrains only almost
everywhere.  Both of those nodes therefore carry, as part of their definition,
the requirement that the relevant integrands be Lebesgue integrable before any
identity or inequality between integrals is asserted; without it the identity of
(5.11) of the thesis and the inequality of clause (4) could both hold by the junk
value that Lebesgue integration assigns to a non-integrable function, and the
predicate would be satisfiable for data for which (C7) is false.  This result is
used as an
explicit hypothesis in Theorem~Theorem~7 of the thesis and is never introduced as an axiom.
\end{definition}
